\documentclass[11pt, twocolumn, landscape, a4paper]{article}

\usepackage[margin=1cm]{geometry}
\usepackage[utf8]{inputenc}
\usepackage[spanish]{babel}
\usepackage{amssymb, amsmath, amsbsy}
\usepackage{tasks}
\usepackage{enumerate}
\usepackage{enumitem}
\usepackage{graphicx}

\usepackage{tikz}
\usetikzlibrary{angles,quotes}

\usepackage[pdftex]{hyperref}
\hypersetup{pdftitle={Trigonometría},pdfauthor={Luis Carrillo Gutiérrez}}

\fontfamily{cmss}\selectfont
\renewcommand{\familydefault}{\sfdefault}
\pagestyle{empty}

\begin{document}
\subsection*{Relación entre lados y ángulos para cualquier triángulo}

\noindent Los siguientes resultados se cumplen para cualquier triángulo de lados \scalebox{1.2}{$a$}, \scalebox{1.2}{$b$}, \scalebox{1.2}{$c$} son ángulos \scalebox{1.1}{$A$}, \scalebox{1.1}{$B$} y \scalebox{1.1}{$C$}. \\

\begin{tikzpicture}
	\coordinate (A) at (0, 0);
	\coordinate (B) at (1.5, 3.25);
	\coordinate (C) at (3.5, 0);
	\draw [thick](B) -- (A) -- (C) -- cycle;

	\draw (-.25, -.2) node{$A$};
	\draw (1.5, 3.56) node{$B$};
	\draw (3.77, -.15) node{$C$};

	\draw pic[draw, angle radius=.35cm] {angle=B--C--A};
	\draw pic[draw, angle radius=.5cm] {angle=A--B--C};
	\draw pic[draw, angle radius=.65cm] {angle=C--A--B};

	\draw (2.88, 1.5) node{$a$};
	\draw (1.8, -.25) node{$b$};
	\draw (.45, 1.6) node{$c$};
\end{tikzpicture}

\begin{itemize}
\item{Ley de los senos
$$
\dfrac{a}{\sen A} = \dfrac{b}{\sen B} = \dfrac{c}{\sen C}
$$
}
\item{Ley de los cosenos
$$
a^{2} = b^{2} + c^{2} - 2 \cdot (b \cdot c) \cdot \cos A
$$
}
\item{Ley de los tangentes
$$
\dfrac{a + b}{a - b} = \dfrac{\tan\; \dfrac1{2}(A + B)}{\tan\; \dfrac1{2}(A - B)}
$$
}
\end{itemize}

% De manera análoga para los otros lados.
\end{document}